\scanpage{57}
% archive sheet number: 56
La \uline{théorie de} \uline{Brauer}-\uline{\unsure{Sévéri|Séveri}}-\uline{\unsure{Châtelet|Chatelet}} \uline{globale}.

1. \unsure{Définition} d'un schéma de Brauer-Sévéri sur $Y$ : \deleted{\unsure{isomorphe}} $\unsure{S|\mathscr{S}}$-isomorphe à $\mathbb{P}^{r}_{Y}$. \inserted{i.e. c'est un schéma de Châtelet} [\unsure{Il est donc} propre, plat sur $Y$, ses fibres sont des Brauer-Sévéri]. Possibilité de la $S$-descente, propreté. Local-isotrivialité.

2. Fibré projectif (\uline{associé} à un faisceau loc. libre) $=$ schéma de Châtelet sur $Y$, \inserted{c'est un Br.-Sé. loc. \unsure{tr.}} muni d'un $\underline{L}$ de rang $1$ qui sur chaque fibre induit \ldots\ (base noeth.); \uline{Corollaire 1} Si $\unsure{Y|S}=\operatorname{Spec}(\unsure{O|0})$, $O$ local complet, alors Châtelet $=$ Brauer-Sévéri. \uline{Corollaire 2} Si $Y$ est non singulier, alors Br. Sévéri localement trivial $=$ \deleted{fibré} projectif $=$ Châtelet qui \deleted{\unsure{un faisceau} \unsure{e.g.} \unsure{en un} \unsure{pt} (de $Y$ \unsure{supposé connexe})} \inserted{est loc. trivial en un pt de la base}
\margin{\unsure{Construction} des diviseurs ``\unsure{linéaires}''}
\inserted{\uline{Cor. 3} S'il existe un diviseur loc. principal \unsure{relatif} comme fibré associé \deleted{\unsure{déjà}} \unsure{dans} (\unsure{sur} $\unsure{\overline{X}=X|X=X'}$).}

3. Fibré dual, \unsure{interprétation} de ses sections comme diviseurs loc. principaux sur $X$ qui sont ``\unsure{spéciaux}''. \uline{Cor. 1} Si $X$ admet une section, il est un fibré projectif [$Y$ noeth.\ldots] De façon plus précise, un $X$ \inserted{de Br.-Sév.,} avec une section, c'est \unsure{la} $=$ chose que la donnée d'un \deleted{fibré} faisceau loc. libre de rang $r$, et d'un sous faisceau loc. libre de rang $r-1$, loc facteur direct\ldots

\uline{Cor. 2} \deleted{\unsure{Autre}} Critère de \unsure{locle} trivialité : existence de sections locales\ldots

4. \uline{Th} Soit $X$ propre et plat sur $Y$. \deleted{\unsure{si}} si $f^{-1}(y)$ est un Br.-Sév., alors il existe un voisinage ouvert $U$ de $y$ t.q. $X|U$ soit Br. Sév.

\uline{Cor.} Châtelet $=$ Brauer-Sévéri.

On \unsure{prouve} d'abord qu'il y a un $U$ t.q. $X|U$ soit Châtelet, puis on \deleted{\unsure{considère}} \unsure{prouve} que Châtelet $\Rightarrow$ Br. Sév., grâce à la considération des $\operatorname{Is}_{Y}(\mathbb{P}^{r}_{Y},X)$.

\{ Relations avec les faisceaux d'algèbres simples centrales\ldots\ Groupe de Brauer global. [\unsure{Consulter} le cours de Serre\ldots] \}

\uline{3 bis} Fibrés des $n$-formes associé à $\deleted{\unsure{Y}}\unsure{X}$, ses sections sur $Y$ sont les diviseurs loc. principaux sur $X$ dont la trace sur \deleted{\unsure{toute}} fibre (ou une fibre, si $Y$ est connexe) est de degré $n$.
